\chapter{Applying Mathematical Thinking to Requirement Specification}
\label{ch:requirement-specification}

This chapter operationalizes the mindset introduced in Chapter~\ref{ch:mathematical-rigor}. We translate requirements into logical propositions, sets, functions, constraints, and graphs so that ambiguity is removed and analysis becomes possible.

\section{Requirements as Logical Propositions}
\label{sec:logical-propositions}

Logic provides building blocks for precise specifications.

\subsection{Basic Logical Operators}
Construct complex requirements from atomic statements using AND, OR, NOT, IMPLIES, and IFF. Truth tables expose contradictions; De Morgan's laws help simplify nested conditions. For example, replacing ``NOT (A OR B)'' with ``NOT A AND NOT B'' clarifies validation rules.

\subsection{Universal and Existential Quantifiers}
Quantifiers ("for all" and "there exists") allow us to describe requirements over collections. ``For all transactions t, there exists an approver a such that a.role = Finance'' encodes compliance checks succinctly.

\subsection{Logical Specification Examples}
Permissions can be expressed as ``User $u$ may view report $r$ iff $u.\text{role} \in \{\text{Analyst}, \text{Admin}\}$ and $r.\text{classification} \le u.\text{clearance}$.'' Data validation might state ``For all records, if currency = 'USD' then amount has two decimal places.''

\subsection{Finding Contradictions}
SAT solvers or theorem provers can check whether a set of logical statements is internally consistent. Even without automation, translating requirements into logic exposes conflicting assumptions quickly.

\section{Set-Based Requirement Modeling}
\label{sec:set-based}

Sets model categories, segments, and capability groupings.

\subsection{Defining User Segments and Personas}
Represent user cohorts as sets (e.g., `PremiumUsers`, `TrialUsers`). Overlaps highlight dual personas and inform entitlement logic.

\subsection{Feature Sets and Dependencies}
Treat features as sets of capabilities. Compatibility constraints become set operations: ``Feature F requires capability set C'' or ``Features X and Y are disjoint.''

\subsection{Data Domains and Constraints}
Domain values such as ISO country codes or allowed statuses are sets. Intersections and unions model multi-criteria filters, and complement sets define invalid inputs.

\subsection{Example: Multi-Criteria Filtering}
An e-commerce search requirement might specify ``$Results = \text{CategorySet} \cap \text{PriceRangeSet} \cap \text{AvailabilitySet}$,'' making the logic explicit and testable.

\section{Functional Specifications}
\label{sec:functional-specs}

Functions map inputs to outputs, clarifying transformations.

\subsection{Pure Functions and Side Effects}
Whenever possible, describe operations as pure functions to simplify reasoning. When side effects exist, document them explicitly (e.g., cache invalidations).

\subsection{Function Signatures and Types}
Specify input and output types. ``$score: \text{UserContext} \times \text{CatalogItem} \rightarrow [0,1]$'' informs implementers and enables compile-time checks.

\subsection{Function Composition}
Complex workflows can be expressed as compositions $(f \circ g)$. Visualizing pipelines as composed functions exposes opportunities for reuse and error handling boundaries.

\subsection{Example: Data Transformation Pipeline}
An ETL requirement might state, ``cleaned = normalize(transform(filter(raw)))'' accompanied by contracts for each function, enabling modular testing.

\section{Constraint-Based Specifications}
\label{sec:constraints}

Declarative constraints state what must hold true without prescribing implementation.

\subsection{Invariants: What Must Always Hold}
Document invariants such as ``$\text{inventory\_available} \ge 0$'' or ``account balance $= \text{credits} - \text{debits}$.'' Monitoring systems can enforce invariants at runtime.

\subsection{Preconditions and Postconditions}
Design by contract ensures functions only execute under valid conditions. For example, ``cancelOrder(order) requires order.state = Pending; ensures order.state = Cancelled and inventory restored.''

\subsection{Temporal Constraints}
Temporal logic captures ordering: ``Payment must be authorized before shipment'' or ``Token expires within 15 minutes of issuance.'' Operators like ``eventually'' or ``always'' help articulate service-level objectives.

\subsection{Example: Reservation System}
Requirements can specify: no overlapping reservations for the same resource, capacity per time slot $\le$ limit, modifications allowed until T-24 hours. Expressing these formally simplifies validation.

\section{State-Based Specifications}
\label{sec:state-based}

State machines capture lifecycle behavior.

\subsection{Finite State Machines}
Define states (Pending, Confirmed, etc.), transitions, and guards. Diagrams reveal unreachable states or missing transitions.

\subsection{State Invariants}
Each state can have invariants, such as ``When an order is Shipped, $\text{trackingId} \neq \text{null}$.'' Automated tests can verify invariants on every transition.

\subsection{Event-Driven Specifications}
Events trigger transitions. Recording guards (conditions) and actions (side effects) ensures clarity about when transitions fire and what they affect.

\subsection{Example: Order Processing Workflow}
Document valid transitions (Pending $\rightarrow$ Confirmed, Confirmed $\rightarrow$ Shipped) and forbidden ones (Delivered $\rightarrow$ Pending). Edge cases like partial shipments can be modeled as substates.

\section{Probabilistic and Stochastic Specifications}
\label{sec:probabilistic}

Some systems require probabilistic descriptions.

\subsection{Probability Distributions}
Specify the distribution governing a variable (e.g., inter-arrival times follow $\text{Poisson}(\lambda)$). These assumptions feed capacity planning and simulation.

\subsection{Expected Values and Variance}
Requirements may target expected outcomes, such as ``Expected customer lifetime value $\ge \$500$ with variance $\le \$50^2$.'' 

\subsection{Markov Chains}
Model user journeys or system states with Markov chains, capturing transition probabilities. Steady-state analysis can estimate churn or queue lengths.

\subsection{Example: Recommendation System Behavior}
Define exploration probabilities, expected click-through per slot, and acceptable variance. Requirements might specify ``Exploration probability decays from 20\% to 5\% over 10 sessions while maintaining expected engagement $\ge$ baseline.''

\section{Graph-Based Specifications}
\label{sec:graph-based}

Graphs represent relationships and dependencies succinctly.

\subsection{Dependency Graphs}
Nodes represent features or services; edges represent dependencies. Requirements can include ``Graph must remain acyclic'' or ``Critical path latency $\le 500\,\text{ms}$.'' 

\subsection{User Journey Graphs}
Model user flows as directed graphs. Bottlenecks or drop-off points become graph metrics (e.g., probability of reaching checkout). Requirements can target improvements on specific edges.

\subsection{Knowledge Graphs}
Domains with rich relationships benefit from ontologies. Requirements might define permissible relationships and cardinalities, ensuring data integrity.

\subsection{Example: Microservices Architecture}
Mapping services and API calls reveals blast radius for failures. Requirements can state ``No service should have more than N synchronous downstream dependencies.''

\section{Formal Specification Languages}
\label{sec:formal-languages}

Dedicated languages provide tooling for rigorous specs.

\subsection{Alloy}
Alloy uses relational logic with automated analyzers that generate counterexamples. Teams can model data models or access rules and validate properties quickly.

\subsection{TLA+}
TLA+ excels at describing concurrent or distributed systems. Model checking uncovers race conditions, deadlocks, or order violations before coding.

\subsection{Z Notation}
Z provides a mathematical schema language ideal for safety-critical domains. Though heavyweight, selected concepts (schemas, invariants) can inspire lighter documentation.

\subsection{Behavior-Driven Development (BDD)}
BDD tools like Cucumber use structured language to create executable specs. They bridge business and technical stakeholders while remaining testable.

\section{Practical Workflow}
\label{sec:practical-workflow}

Incorporate formalization incrementally.

\subsection{Collaborative Specification Sessions}
Adopt the ``Three Amigos'' practice: PM, engineer, and tester co-author requirements, translating user stories into logic, tables, or diagrams together.

\subsection{Iterative Refinement}
Start with informal prose and formalize areas where ambiguity persists. Revisit specs as experiments reveal new information.

\subsection{Tool-Assisted Specification}
Use modeling tools, linters, or IDE plugins that highlight inconsistencies. Integrate specs with issue trackers so updates stay synchronized.

\subsection{Living Documentation}
Treat specifications as code: version-control them, review via pull requests, and run automated checks in CI. This keeps docs aligned with implementation.

\section{Industry Scenarios}
\label{sec:req-spec-examples}

\paragraph{Fintech.} A banking team represented entitlements as set expressions (PremiumUsers $\cap$ SecuritiesCleared) and logical guards, preventing unauthorized trading flows; audits passed smoothly. A rival project codified access rules in prose only, leading to inconsistent implementations across squads and a regulatory finding.

\paragraph{Healthcare.} A clinical data platform defined data lineage as functions mapping source $\rightarrow$ transformation $\rightarrow$ patient record, plus invariants (no PHI leaves region). This clarity accelerated compliance reviews. A patient-portal rewrite skipped formal constraints, and duplicate record merges corrupted histories until predicates were added retroactively.

\paragraph{E-commerce.} A logistics system captured fulfillment states in state machines and sequence diagrams, making it easy to spot missing transitions for split shipments. Another retailer relied on narrative specs for promotions; edge cases like overlapping discounts caused financial leakage until decision tables were introduced.

\section{Summary}
\label{sec:req-spec-summary}

Mathematical notation turns requirements into analyzable artifacts. Logic, sets, functions, constraints, state machines, and graphs each offer lenses to eliminate ambiguity. Choose the level of formality appropriate for the risk profile, and keep specifications living alongside the codebase.

\subsection{Action Checklist}
\begin{enumerate}
    \item Select a complex requirement and translate it into logical statements or set relationships.
    \item Build or update a state diagram or sequence diagram for the riskiest workflow.
    \item Ensure specs live in version control with review checkpoints similar to code changes.
\end{enumerate}

\paragraph{Try this now (5 min).} Open a current user story and underline every conditional phrase. Convert them into a small truth table or bullet list of cases to verify coverage.

\section{Day in the Life: Maya Hosts a Spec Jam}
Midweek, Maya realizes the new triage workflow spec lives in three conflicting docs. She schedules a ``spec jam'' with an engineer and QA lead. Together they project the spec into their repo and translate the prose into logical statements—``IF vital\_risk $\ge$ 0.8 AND consent\_captured THEN route to tele-cardiology.'' QA notes a missing scenario (no consent), so they add an exception and assign telemetry requirements. Maya then opens a pull request for the spec, tags compliance, and links automated tests to the predicates. The shared, version-controlled artifact becomes the single truth source, and new engineers ramp faster because the logic is explicit.

\section{Visual Guide}
\label{sec:reqspec-visuals}
\begin{itemize}
    \item \textbf{Spec-to-Code Pipeline}—diagram how a requirement flows from collaborative spec session to version control to automated checks.
    \item \textbf{Logic Table Example}—small truth table derived from Maya’s triage rule, showing coverage of all cases.
    \item \textbf{Set/Function Diagram}—visualizing user segments or data flows to illustrate how set theory clarifies entitlements.
\end{itemize}

\section{Interactive Elements}
\label{sec:reqspec-interactive}
\begin{itemize}
    \item \textbf{Self-Assessment}—checklist evaluating current specs for clarity, traceability, and test linkage.
    \item \textbf{Reflection Prompt}—``Which spec lacks ownership?'' Encourage writing down responsible reviewers and missing approvals.
    \item \textbf{Pause and Practice}—guided exercise for running a ``spec jam'' session with prompts for logic translation.
    \item \textbf{Implementation Planner}—template for introducing docs-as-code workflows (repo setup, review cadence, tooling).
\end{itemize}
\section{Industry Adaptations}
\label{sec:reqspec-industry}
\paragraph{Regulated Industries.} Specifications must capture traceability—link requirements to regulations and ensure logic is reviewable by auditors. Use requirement IDs consistently and maintain change logs with approvals.

\paragraph{Startups vs. Enterprises.} Startups can hold lightweight ``three amigos'' sessions and store specs in the repo to keep alignment high. Enterprises should institutionalize spec reviews with sign-offs from security/legal where appropriate.

\paragraph{B2B vs. B2C.} B2B specs should include integration contracts, SLAs, and partner dependencies. B2C specs should anticipate scale and localization—explicitly document segments and user journeys.

\paragraph{Hardware-Software Integration.} Maintain joint diagrams showing signal flow between hardware components and software services. Include hardware specs (pinouts, timing) alongside API contracts to prevent mismatches.
\section{Warning Signs and Recovery}
\label{sec:reqspec-warning}
\begin{description}
    \item[\textbf{Warning sign}] Specification reviews focus on style or formatting while contradictions between requirements slip through.
    \item[\textbf{Recovery}] Introduce a checklist that explicitly verifies logic consistency, completeness of states, and ownership of invariants.
    \item[\textbf{Warning sign}] Documentation drifts away from implementation because specs live in static documents.
    \item[\textbf{Recovery}] Manage specs like code—version control them, require pull-request reviews, and link them to automated checks where possible.
\end{description}

\section{Triple-Lens Takeaways}
\label{sec:req-spec-triple}

\begin{description}
    \item[\textbf{Busy PM}] Reuse the collaborative specification workflow to schedule short ``three amigos'' sessions—translate sticky requirements into logic or diagrams only where ambiguity persists.
    \item[\textbf{Technical Lead}] Treat each modeling technique as part of a toolkit: map auth rules to logic, data entitlements to sets, workflows to state machines, and wire them into tests or linters.
    \item[\textbf{Executive}] Ask teams to show the living specs for the riskiest initiatives; the presence of version-controlled, reviewable artifacts is a proxy for operational maturity and risk control.
\end{description}
